\documentclass[a4paper,11pt]{ltjsarticle}
\usepackage{luatexja}
\usepackage{luatexja-fontspec}
\usepackage{amsmath,amssymb,amsthm}
\usepackage{mathtools}
\usepackage{booktabs}
\usepackage{longtable}
\usepackage{array}
\usepackage{hyperref}
\usepackage{graphicx}
\usepackage{float}
\usepackage{geometry}
\geometry{margin=25mm}

\hypersetup{
  colorlinks=true,
  linkcolor=blue,
  urlcolor=blue,
  citecolor=blue
}

\theoremstyle{definition}
\newtheorem{theorem}{定理}[section]
\newtheorem{proposition}[theorem]{命題}
\newtheorem{lemma}[theorem]{補題}
\newtheorem{corollary}[theorem]{系}
\newtheorem{definition}[theorem]{定義}
\newtheorem{remark}[theorem]{注意}
\newtheorem{observation}[theorem]{観察}

\setlength{\parindent}{1em}
\setlength{\parskip}{0.5em}

\providecommand{\tightlist}{\setlength{\itemsep}{0pt}\setlength{\parskip}{0pt}}
\begin{document}

\section{中心投影における曲率半径の極限についての考察}\label{ux4e2dux5fc3ux6295ux5f71ux306bux304aux3051ux308bux66f2ux7387ux534aux5f84ux306eux6975ux9650ux306bux3064ux3044ux3066ux306eux8003ux5bdf}

\textbf{\(R \to \infty\) と \(R \to 0\) の場合}

\begin{center}\rule{0.5\linewidth}{0.5pt}\end{center}

\textbf{著者}: 木原 範昭 (WF System Co., Ltd.)

\textbf{日付}: 2026年4月

\textbf{DOI：} \href{https://doi.org/10.5281/zenodo.19526549}{10.5281/zenodo.19526549}

\begin{center}\rule{0.5\linewidth}{0.5pt}\end{center}

\subsection{要旨}\label{ux8981ux65e8}

本稿では、先行論文 {[}1{]}, {[}2{]}, {[}3{]} で確立された中心投影フレームワークにおいて、曲率半径 \(R\) の極限 \(R \to \infty\) および \(R \to 0\) の両極限について考察する。両極限において、誘導計量、曲率テンソル、および導出される Einstein 方程式の各項は幾何学的に正則である。一方、\(R \to 0\) の極限については、論文 {[}2{]} で示された測地線偏差の上限との関係において、構造的な緊張が生じる。この緊張を幾何学の枠組みの内部で解消する一つの方法として、\textbf{主観空間の入れ子構造}が先行論文 {[}1{]}, {[}2{]} の既存の幾何学的自由度から直接構成可能であることを示す。本稿で述べるのは、このような構造の幾何学的存在可能性のみであり、その物理的実在性については何も主張しない。

\begin{center}\rule{0.5\linewidth}{0.5pt}\end{center}

\subsection{1. 導入}\label{ux5c0eux5165}

先行論文 {[}1{]} では、\((n+1)\) 次元ユークリッド背景空間 \(\mathbb{R}^{n+1}\) における半径 \(R\) の超球面 \(S^n(R)\) と接超平面 \(\Pi_R\) の間の中心投影

\[
\Phi : \Pi_R \longrightarrow S^n(R)
\]

を定義し、\(\Pi_R\) 上に誘導される計量

\[
g_{\mu\nu}(x) = \frac{R^2}{l^2}\left(\delta_{\mu\nu} - \frac{x_\mu x_\nu}{l^2}\right), \qquad l^2 = R^2 + |x|^2
\]

から、Einstein 方程式

\[
G_{\mu\nu} + \Lambda\, g_{\mu\nu} = 0, \qquad \Lambda = \frac{3}{R^2}
\]

が導出されることを示した。

論文 {[}2{]} では、このフレームワークが持つ対称性(離散安定性、軸の対等性、測地線偏差の構造、主観座標系の変換可能性)を証明し、特に主観空間内の観測者が測定する測地線偏差 \(|\xi|^2 = g_{\mu\nu}\xi^\mu \xi^\nu\) に曲率半径 \(R\) に依存する上限があることを示した。

論文 {[}3{]} では、複数の主観空間が同一の背景空間を共有する場合の観測の相対性を議論した。

これらの論文において、曲率半径 \(R\) は自由パラメータとして残されており、\(R \in (0,\infty)\) の全域で構成は well-defined である。本稿の目的は、両端の極限

\[
R \to \infty, \qquad R \to 0
\]

における幾何学的構造を検証し、そこから自然に浮上する問いを記録することにある。

\begin{center}\rule{0.5\linewidth}{0.5pt}\end{center}

\subsection{\texorpdfstring{2. 極限 \(R \to \infty\) の場合}{2. 極限 R \textbackslash to \textbackslash infty の場合}}\label{ux6975ux9650-r-to-infty-ux306eux5834ux5408}

\subsubsection{2.1 幾何学的正則性}\label{ux5e7eux4f55ux5b66ux7684ux6b63ux5247ux6027}

\(R \to \infty\) の極限において、誘導計量 \(g_{\mu\nu}\) は形式的に

\[
g_{\mu\nu}(x) \longrightarrow \delta_{\mu\nu}
\]

に漸近し、主観空間 \(\Pi_R\) は平坦ユークリッド空間に近づく。宇宙項は

\[
\Lambda = \frac{3}{R^2} \longrightarrow 0
\]

となり、Einstein 方程式の各項は数学的に well-defined な漸近を示す。曲率テンソルの各成分は連続的にゼロに近づき、論文 {[}1{]} の幾何学的構成はこの極限で破綻しない。

\subsubsection{2.2 幾何学の外側に残る問い}\label{ux5e7eux4f55ux5b66ux306eux5916ux5074ux306bux6b8bux308bux554fux3044}

しかし、本構成の物理的解釈を試みる場合、以下の問いが立つ。これらは幾何学的フレームワークの外側の問いであり、本稿では扱わない。

\textbf{問い 2.1}　背景空間 \(\mathbb{R}^{n+1}\) はユークリッド符号を持つ。\(R \to \infty\) の極限で主観空間が平坦化するとき、時間に対応する軸が構造的に区別されるかどうかは、純粋幾何学の枠組みでは決定されない。

\textbf{問い 2.2}　\(\Lambda \to 0\) は滑らかな漸近であるが、主観空間の体積は \(R \to \infty\) で発散する。これらの量の積として現れうる量(真空エネルギーと解釈されうる量)の振る舞いは、幾何学の外側で定義される解釈に依存する。

\textbf{問い 2.3}　\(\Lambda\) を中心投影とは異なるフレームワークから導出する場合、その振る舞いが \(R \to \infty\) の極限で \(3/R^2\) と整合するかどうかは、各フレームワークの構造に依存する。

\textbf{要点}: 幾何学的には正則であるが、物理解釈では発散する可能性が否定できない。

\begin{center}\rule{0.5\linewidth}{0.5pt}\end{center}

\subsection{\texorpdfstring{3. 極限 \(R \to 0\) の場合}{3. 極限 R \textbackslash to 0 の場合}}\label{ux6975ux9650-r-to-0-ux306eux5834ux5408}

\subsubsection{3.1 幾何学的正則性}\label{ux5e7eux4f55ux5b66ux7684ux6b63ux5247ux6027-1}

\(R \to 0\) の極限において、\(\Lambda = 3/R^2\) は発散する。しかし、誘導計量 \(g_{\mu\nu}\)、曲率テンソル、および Einstein 方程式の各項は、\(R > 0\) の任意の値で well-defined であり、\(R \to 0\) の極限も幾何学的に正則な漸近として定義される。論文 {[}1{]} の構成はこの極限でも破綻しない。

\subsubsection{3.2 測地線偏差の上限との構造的緊張}\label{ux6e2cux5730ux7ddaux504fux5deeux306eux4e0aux9650ux3068ux306eux69cbux9020ux7684ux7dcaux5f35}

論文 {[}2{]} において、主観空間内の観測者が測定する測地線偏差 \(|\xi|\) には曲率半径 \(R\) に依存する上限が存在することが示された。すなわち、観測者が測定可能な量は主観空間の大きさのオーダーによって制限される。

\(R \to 0\) の極限において、この上限はゼロに漸近する一方、局所曲率は \(1/R^2\) のオーダーで発散する。したがって、主観空間内の観測者が測定可能な量と、その観測者が住む主観空間の局所幾何学的構造の間に、構造的な緊張が生じる。

\subsubsection{3.3 幾何学の外側に残る問い}\label{ux5e7eux4f55ux5b66ux306eux5916ux5074ux306bux6b8bux308bux554fux3044-1}

\textbf{問い 3.1}　\(\Lambda = 3/R^2 \to \infty\) は幾何学的な量の発散として定義されるが、この発散を物理的にどう解釈するかは幾何学の外側の問いである。

\textbf{問い 3.2}　主観空間が点に縮む極限において、観測者が主観空間の内部に存在するという構造そのものが保たれるかどうかは、幾何学の内部では決定されない。

\textbf{要点}: 幾何学的には正則であるが、物理解釈では発散する可能性が否定できず、かつ §3.2 の緊張が観測者の構造との整合性について問いを残す。

\begin{center}\rule{0.5\linewidth}{0.5pt}\end{center}

\subsection{4. 緊張の解消: 主観空間の入れ子構造}\label{ux7dcaux5f35ux306eux89e3ux6d88-ux4e3bux89b3ux7a7aux9593ux306eux5165ux308cux5b50ux69cbux9020}

§3.2 で指摘した緊張を、幾何学の枠組みの内部で解消する一つの方法を示す。本節で述べるのは、先行論文 {[}1{]}, {[}2{]} の既存の定理の直接的な帰結であり、新たな公理の導入を要しない。

\subsubsection{4.1 入れ子構造の構成}\label{ux5165ux308cux5b50ux69cbux9020ux306eux69cbux6210}

論文 {[}1{]} の中心投影フレームワークにおいて、主観空間 \(\Pi_R\) は軸 \(e\) の選択によって定義される。論文 {[}2{]} 定理 3.1(軸の対等性)により、この軸の選択は幾何学的に一意でなく、任意の軸 \(e'\) を新たな投影中心として選ぶ自由度がある。

この自由度は、主観空間の内部の任意の点に対しても適用可能である。すなわち、\(\Pi_R\) 内の任意の点 \(p\) に対し、\(p\) を新たな軸の立ち上がり点とする別の中心投影

\[
\Phi' : \Pi_{R'} \longrightarrow S^n(R')
\]

を定義することができる。ここで新しい曲率半径 \(R' > 0\) は任意に選べる。新しい主観空間 \(\Pi_{R'}\) は、論文 {[}1{]} の構成と同じ幾何学的構造(誘導計量、曲率テンソル、Einstein 方程式)を持つ。

\subsubsection{4.2 存在命題}\label{ux5b58ux5728ux547dux984c}

\textbf{命題 4.1}(入れ子構造の存在)　論文 {[}1{]} で定義された中心投影フレームワークにおいて、任意の主観空間 \(\Pi_R\) と、その内部の任意の点 \(p \in \Pi_R\) に対し、次の性質を満たす主観空間 \(\Pi_{R'}\) が任意の \(R' > 0\) に対して存在する:

\begin{enumerate}
\def\labelenumi{\arabic{enumi}.}
\tightlist
\item
  \(\Pi_{R'}\) は \(p\) を中心とする新たな中心投影によって定義される。
\item
  \(\Pi_{R'}\) は論文 {[}1{]} で定義された構成と同じ幾何学的構造を持つ。
\item
  \(\Pi_{R'}\) は \(\Pi_R\) と同じ背景空間 \(\mathbb{R}^{n+1}\) から構成される。
\end{enumerate}

\textbf{証明の概略}　論文 {[}2{]} 定理 3.1 により、\(\Pi_R\) の任意の内部点 \(p\) において、\(p\) を立ち上がり点とする軸 \(e'\) を選ぶことができる。この \(e'\) から、論文 {[}1{]} と同じ手続きで半径 \(R'\) の超球面 \(S^n(R')\) および接超平面 \(\Pi_{R'}\) を構成できる。構成は論文 {[}1{]} の定理によって well-defined である。\(\square\)

\subsubsection{4.3 入れ子の反復可能性}\label{ux5165ux308cux5b50ux306eux53cdux5fa9ux53efux80fdux6027}

命題 4.1 の構成は反復可能である。\(\Pi_{R'}\) の内部点 \(p'\) において、さらに \(\Pi_{R''}\) を立ち上げることができ、以下同様である。この反復は論文 {[}1{]}, {[}2{]} のフレームワークの内部で閉じており、追加の公理を必要としない。

\subsubsection{4.4 §3.2 の緊張との関係}\label{ux306eux7dcaux5f35ux3068ux306eux95a2ux4fc2}

命題 4.1 が存在することにより、§3.2 で指摘した緊張は次のように幾何学の内部で整理される:

主観空間 \(\Pi_R\) の \(R \to 0\) の極限的な領域に対して、その領域を新たな主観空間 \(\Pi_{R'}\) の立ち上がり点として解釈する自由度が幾何学的に存在する。親の主観空間 \(\Pi_R\) から見れば縮退する領域が、新たな中心投影 \(\Pi_{R'}\) の内部から見れば well-defined な幾何学的構造を持つ。

この整理によって、親の主観空間内で収容困難となる構造的要素は、子の主観空間に対する投影として記述できる幾何学的自由度を持つ。

\begin{center}\rule{0.5\linewidth}{0.5pt}\end{center}

\subsection{5. 考察}\label{ux8003ux5bdf}

本稿では、先行論文 {[}1{]}, {[}2{]}, {[}3{]} の中心投影フレームワークの曲率半径の両極限を検証し、§4 において主観空間の入れ子構造が既存のフレームワークの幾何学的自由度として構成可能であることを示した。

本稿の主張は以下に限定される:

\begin{enumerate}
\def\labelenumi{\arabic{enumi}.}
\tightlist
\item
  両極限において幾何学的構成は破綻しない。
\item
  \(R \to 0\) の極限において §3.2 の構造的緊張が生じる。
\item
  この緊張を幾何学の内部で整理する方法として、§4 の入れ子構造の存在命題が成立する。
\end{enumerate}

本稿は、このような入れ子構造が物理的に実在するとは主張しない。§2.2 および §3.3 で提示された問いについても、物理的解釈は幾何学の外側の問題であり、本稿では扱わない。

これらの問いへの解答、および入れ子構造の物理的意味付けは、読者の判断に委ねられる。

\begin{center}\rule{0.5\linewidth}{0.5pt}\end{center}

\subsection{6. 結論}\label{ux7d50ux8ad6}

先行論文 {[}1{]}, {[}2{]}, {[}3{]} の中心投影フレームワークにおいて、曲率半径 \(R\) の両極限 \(R \to \infty\), \(R \to 0\) は幾何学的に正則である。\(R \to 0\) の極限で生じる測地線偏差の上限との構造的緊張は、同フレームワークの既存の自由度から構成される主観空間の入れ子構造によって、幾何学の内部で整理可能である。本構成の物理的解釈は本稿の対象外であり、読者に委ねる。

\begin{center}\rule{0.5\linewidth}{0.5pt}\end{center}

\subsection{参考文献}\label{ux53c2ux8003ux6587ux732e}

{[}1{]} 木原 範昭,「中心投影による幾何学的定式化」, Zenodo, DOI: 10.5281/zenodo.19427780 (2025).

{[}2{]} 木原 範昭,「中心投影フレームワークにおける対称性」, Zenodo, DOI: 10.5281/zenodo.19434932 (2025).

{[}3{]} 木原 範昭,「複数主観空間における観測の相対性」, Zenodo, DOI: 10.5281/zenodo.19435162 (2025).

\begin{center}\rule{0.5\linewidth}{0.5pt}\end{center}

\end{document}
